% ============================================================================
% PREAMBULE
% Fichier : preambule/preambule.tex
% ============================================================================

\chapter*{PR\'EAMBULE}
\markboth{Pr\'eambule}{}
\addcontentsline{toc}{chapter}{Pr\'eambule}
\thispagestyle{preliminaire}

\begin{flushright}
  \begin{minipage}{0.72\textwidth}
    \itshape
    \guillemotleft~La cryptographie n'existe pas dans le vide. Elle est au service
    d'institutions humaines qui op\`erent sous des contraintes cognitives,
    organisationnelles et proc\'edurales.~\guillemotright
  \end{minipage}
\end{flushright}
\vspace{0.3cm}
\begin{flushright}
  \small\textit{- Principe fondateur de la Cryptographie
  Institutionnelle~\cite{MINKA2026CLO}}
\end{flushright}

\vspace{1.0cm}

Un arbitre institutionnel, qu'il soit juge, auditeur ou r\'egulateur, se trouve
confront\'e \`a une preuve fond\'ee sur des principes cryptographiques.
L'algorithme a effectu\'e toutes les v\'erifications n\'ecessaires et la
s\'ecurit\'e du sch\'ema n'est pas en cause. L'arbitre se trouve pourtant
dans l'incapacit\'e de rendre sa d\'ecision. Sa comp\'etence n'est pas en
cause. La preuve, bien que math\'ematiquement irr\'eprochable, reste opaque
du point de vue institutionnel. Elle atteste de quelque chose sans en pr\'eciser la nature,
dans les termes requis pour pouvoir statuer. Ce d\'ecalage entre la
rigueur du protocole et la clart\'e du verdict constitue le point de
d\'epart de cette th\`ese.

Le terme \guillemotleft~arbitre institutionnel~\guillemotright~fait r\'ef\'erence
\`a ce destinataire final, volontairement g\'en\'erique, qu'il s'agisse d'un
juge civil ou p\'enal, d'un auditeur des finances publiques, d'une commission
\'electorale ou d'un r\'egulateur sectoriel. Cet arbitre n'est ni cryptographe
ni sp\'ecialiste en informatique. Il est celui pour qui la preuve est
\'elabor\'ee. Sa conviction est la base de la d\'ecision qui s'impose. La
cryptographie a longtemps n\'eglig\'e cet aspect. Ce travail vise \`a en faire
une pr\'eoccupation centrale de la discipline.

Ce personnage est pr\'esent tout au long de la th\`ese. Au d\'ebut de chaque
chapitre, une note succincte permet de le repositionner dans le cadre de la
progression canonique
$\mathrm{F}{\to}\mathrm{B}{\to}\mathrm{E}{\to}\mathrm{S}{\to}\mathrm{C}{\to}\mathrm{X}$,
repr\'esentant respectivement formaliser, borner, \'etendre, saturer,
certifier, puis eXp\'erimenter. Ce fil directeur rappelle que la
preuve est con\c{c}ue pour un acteur humain dont la d\'ecision d\'epasse le
simple calcul.

L'arbitre n'est pas l'auteur de cette thèse. Le parcours professionnel décrit à la sous-section~\ref{subsubsec:experience_auteur}
a fourni le terrain d'observation, dont ce personnage est l'abstraction théorique.

\medskip
\noindent\textit{Cette sc\`ene ouvre la th\`ese.}

\vspace{0.8cm}

\begin{boiteArbitre}
L'arbitre franchit le seuil de la salle. Devant lui se pr\'esente une preuve
cryptographique attestant l'int\'egrit\'e de milliers de transactions. Un
algorithme en a valid\'e la v\'eracit\'e en quelques millisecondes. Aucun
adversaire op\'erant en temps polynomial ne saurait la contrefaire. Pourtant,
il lui revient de prononcer un verdict, ce qui suppose de saisir pr\'ecis\'ement ce
que la preuve \'etablit, au-del\`a du seul fait qu'une certification a eu
lieu. Or il en est incapable. Sur le plan
math\'ematique, la preuve est irr\'eprochable, tandis que, du point de vue
institutionnel, elle demeure opaque. Ce paradoxe, qu'il \'eprouve sans
parvenir \`a le formuler, constitue l'objet de cette th\`ese.
\end{boiteArbitre}